% !TeX program = pdflatex
% ═══════════════════════════════════════════════════════════════
%   Ioannis (John) Rizos — résumé for AI expert-evaluation work
%   Deliberately plain: single column, no tables, standard headings,
%   so that automated résumé parsers read it correctly.
% ═══════════════════════════════════════════════════════════════
\documentclass[a4paper,10.5pt]{article}

\usepackage[utf8]{inputenc}
\usepackage[T1]{fontenc}
\usepackage{charter}
\usepackage[protrusion=true,expansion=false]{microtype}
\usepackage[english]{babel}
\usepackage[a4paper,top=1.5cm,bottom=1.5cm,left=1.9cm,right=1.9cm]{geometry}
\usepackage{enumitem}
\usepackage{xcolor}
\usepackage[hidelinks]{hyperref}

\definecolor{rule}{HTML}{555555}

\setlength{\parindent}{0pt}

% Identifying footer: name plus "page N of M" reads cleanly if the
% pages are ever separated, and parses as text rather than a bare digit.
\usepackage{lastpage}
\usepackage{fancyhdr}
\pagestyle{fancy}
\fancyhf{}
\renewcommand{\headrulewidth}{0pt}
\renewcommand{\footrulewidth}{0pt}
\fancyfoot[C]{\footnotesize\color{rule}%
  Ioannis (John) Rizos --- page \thepage\ of \pageref{LastPage}}
\setlength{\footskip}{22pt}

\newcommand{\sect}[1]{%
  \vspace{9pt}%
  {\large\bfseries #1}\par
  \vspace{2pt}{\color{rule}\hrule height 0.6pt}\vspace{6pt}%
}
\newcommand{\jobline}[2]{\textbf{#1} \hfill #2\par\vspace{1pt}}
\setlist[itemize]{leftmargin=1.15em, itemsep=1.5pt, parsep=0pt, topsep=2pt}

\begin{document}

% ── HEADER ───────────────────────────────────────────────
{\huge\bfseries Ioannis (John) Rizos}\par
\vspace{3pt}
Professor of Theoretical Physics, Ph.D. --- Ioannina, Greece\par
\vspace{2pt}
irizos@proton.me \textbullet\ +30 6945297106 \textbullet\
ORCID 0000-0001-7586-1700 \textbullet\
\href{https://inspirehep.net/authors/991609}{INSPIRE-HEP} \textbullet\
\href{https://scholar.google.com/citations?user=7TJ7ZnAAAAAJ}{Google Scholar}

\sect{Summary}

Theoretical physicist with nearly thirty years of university teaching,
examining and research supervision. Extensive practical experience in
exactly the work that expert evaluation requires: authoring original
physics and mathematics problems, writing and recording step-by-step
model solutions, designing assessment items and marking rubrics, and
judging whether a technical argument is correct and where it fails.
Author of 71 peer-reviewed journal articles; referee for
\textit{Nuclear Physics B} and \textit{JHEP} since 2008. Recent research
applies machine learning, genetic algorithms and quantum annealing to
large-scale classification problems.

\sect{Skills}

\textbf{Subject expertise.} Quantum mechanics; quantum field theory;
classical mechanics; electrodynamics; statistical physics; group theory;
string theory and particle phenomenology. Mathematics for machine
learning: linear algebra, calculus, probability, statistics,
optimisation, numerical methods.

\vspace{3pt}
\textbf{Evaluation and assessment.} Original problem authoring; worked
solution writing; multiple-choice item and distractor design; rubric
construction; computer-based assessment systems; peer review of
technical manuscripts; doctoral and master's thesis examination.

\vspace{3pt}
\textbf{Technical.} Python, Mathematica, FORTRAN, \LaTeX;
genetic algorithms, quantum annealing, machine-learning classification;
UNIX/Linux.

\vspace{3pt}
\textbf{Languages.} Greek (native), English (fluent), French (very good).

\sect{Experience}

\jobline{Professor of Theoretical Physics --- University of Ioannina}{1996 -- present}
\textit{Assistant Professor 1996--2004; Associate Professor 2004--2014;
Professor 2014--present}
\begin{itemize}
\item Designed, taught and examined 15 undergraduate and graduate courses
  across quantum theory, classical mechanics, electrodynamics, group
  theory and computational physics --- more than 120 semester-courses in
  total.
\item Authored and marked written examinations and problem sets
  continuously since 1996, including the design of problems
  intended to discriminate between levels of understanding.
\item Supervised 15 master's theses and one doctoral dissertation; served
  on numerous three-member advisory and seven-member doctoral
  examination committees.
\item Created and taught \textit{Computational Mathematics and Physics}
  for 22 semesters, covering numerical and symbolic computation in
  Python and Mathematica.
\end{itemize}

\vspace{5pt}
\jobline{Problem Author and Scientific Editor --- Papazisis Publications}{3rd ed.\ 2019, 4th ed.\ 2022}
\textit{Greek edition of Young \& Freedman,} Sears and Zemansky's
University Physics with Modern Physics \textit{(14th ed.)}
\begin{itemize}
\item One of a twelve-member academic team responsible for the scientific
  editing and adaptation of a two-volume, 1{,}880-page university physics
  textbook (3rd and 4th Greek editions).
\item Wrote original physics problems not present in the American edition,
  and personally created and recorded model video solutions; the team
  produced 254 such solutions, delivered via QR codes printed in the text.
%\item Contributed to a set of 860 teaching slides covering every chapter,
%  published online as course material.
\end{itemize}

\vspace{5pt}
\jobline{Tutor and Module Coordinator --- Hellenic Open University}{2004 -- 2025}
\begin{itemize}
\item \textit{Mathematics for Machine Learning} (DAMA50), M.Sc.\ in Data
  Science and Machine Learning, taught in English, 2021--2025; Module
  Coordinator 2022--24. Linear algebra, calculus, probability, statistics
  and optimisation as the mathematical basis of machine learning.
\item Member of the programme's Curriculum Committee, 2021--2024.
\item \textit{Quantum Physics} (2017--2021) and \textit{Classical
  Physics~II} (2004--2017); Module Coordinator for both. Distance
  education throughout, with written assessment as the primary mode.
\item Authored multiple-choice question banks and produced 24 hours of
  recorded physics teaching material.
\end{itemize}

\vspace{5pt}
\jobline{Technical Review, Evaluation and Examination}{2008 -- present}
\begin{itemize}
\item Referee, on invitation, for \textit{Nuclear Physics B},
  \textit{Journal of High Energy Physics}, \textit{International Journal
  of Modern Physics A}, \textit{Modern Physics Letters A} and
  \textit{Universe} --- assessing the correctness of technical arguments
  and derivations.
\item Editorial boards: \textit{International Journal of Modern Physics A}
  and \textit{Modern Physics Letters A} (2020--), \textit{Particles}
  (2021--).
\item \textbf{Faculty appointment and promotion committees.} Member of
  three-person evaluation committees (\textit{eisigitikes epitropes})
  that assess candidates' full research records and produce the written,
  reasoned report on which appointment decisions are based --- three such
  committees in 2025--26 alone, at the University of Ioannina, the
  Aristotle University of Thessaloniki and the NCSR Demokritos research
  centre. Also serves on the electoral bodies that decide academic
  appointments and promotions.
\end{itemize}

\vspace{5pt}
\jobline{Leadership roles}{2007 -- present}
\begin{itemize}
\item President, Southeastern European Network in Mathematical and
  Theoretical Physics (SEENET-MTP), 2018--present; re-elected 2025.
\item Head, Department of Physics, University of Ioannina, 2013--2017;
  Deputy Head 2009--2013; Chair of the Internal Evaluation Group since
  2022.
\end{itemize}

\sect{Assessment Systems and Education Research}

\begin{itemize}
\item \textbf{TELEDIA} --- co-developed a dynamic computer-based
  assessment system for university physics, in collaboration with the
  Department of Education (published 2002).
\item Published studies on the evolution of physics students' conceptions
  of Newtonian mechanics, and on the use of digital technologies in
  assessment (2002--2007).
\end{itemize}

\sect{Education}

\textbf{Ph.D.\ in Physics} --- University of Ioannina, 1991.
Thesis on the construction and phenomenological analysis of
supersymmetric unified models from four-dimensional superstring theories.\par
\vspace{2pt}
\textbf{B.Sc.\ in Physics} --- University of Ioannina, 1984.\par
\vspace{2pt}
Postdoctoral researcher, École Polytechnique, France (1991--93) and
SISSA/ISAS, Italy (1993--96). Marie Curie Fellow (1995--96, 2003).
Scientific Associate, CERN (2013).

\sect{Selected Publications}

71 articles in peer-reviewed international journals; 5{,}886 citations,
h-index 39 (Google Scholar). Most relevant to computational and
machine-learning work:
\begin{itemize}
\item ``String Model Building on Quantum Annealers,'' S.~A.~Abel,
  L.~A.~Nutricati and J.~Rizos, \textit{Fortschritte der Physik} 71
  (2023) 2300167.
\item ``Towards Machine Learning in the Classification of $Z_2\times Z_2$
  Orbifold Compactifications,'' A.~E.~Faraggi, G.~Harries, B.~Percival
  and J.~Rizos, \textit{J.~Phys.: Conf.~Ser.} 1586 (2020) 012032.
\item ``Genetic Algorithms and the Search for Viable String Vacua,''
  S.~Abel and J.~Rizos, \textit{JHEP} 1408 (2014) 010.
\item ``MERLIN: A Portable System for Multidimensional Minimization,''
  G.~A.~Evagelakis, J.~Rizos, I.~E.~Lagaris and I.~N.~Demetropoulos,
  \textit{Computer Physics Communications} 46 (1987) 401--415.
\end{itemize}

Full list: \href{https://inspirehep.net/authors/991609}{inspirehep.net/authors/991609}

\end{document}
